% ৬. উচ্চস্তরের ভাষা
\section{উচ্চস্তরের ভাষা}

\begin{learningbox}
এই টপিক শেষে শিক্ষার্থী high-level language-এর বৈশিষ্ট্য, সুবিধা ও translator প্রয়োজনীয়তা ব্যাখ্যা করতে পারবে।
\end{learningbox}

\begin{definitionbox}{উচ্চস্তরের ভাষা}
মানুষের ভাষা ও গণিতের notation-এর কাছাকাছি যে programming language, তাকে high-level language বলা হয়। এটি সাধারণত machine independent এবং compiler/interpreter দিয়ে machine code-এ রূপান্তরিত হয়।
\end{definitionbox}

\begin{keypointbox}{সুবিধা}
\begin{itemize}
\item লেখা ও বোঝা সহজ।
\item program development দ্রুত।
\item portability তুলনামূলক বেশি।
\item বড় application তৈরি করা সুবিধাজনক।
\end{itemize}
\end{keypointbox}

\begin{examplebox}{উদাহরণ}
C, C++, Java, Python, BASIC, Fortran, Visual Basic ইত্যাদি high-level language হিসেবে ব্যবহৃত হয়।
\end{examplebox}
